% ===========================================================================
%  Raport — Faza 4: Scalare la Modele Mai Mari
%  ViT-Tiny / ViT-Small / ViT-Base pe ImageNet-1k validation
%  NVIDIA A100 (Google Colab Pro+)
%  Tudose Alexandru · Lucrare de Licență 2026
% ===========================================================================
\documentclass[12pt,a4paper]{article}

% --- Encoding & fonts (XeTeX / tectonic) ---
\usepackage{fontspec}
\defaultfontfeatures{Ligatures=TeX}

% --- Language ---
\usepackage{polyglossia}
\setmainlanguage{romanian}

% --- Page geometry ---
\usepackage[
  top=2.5cm, bottom=2.5cm,
  left=2.5cm, right=2.5cm
]{geometry}

% --- Math ---
\usepackage{amsmath, amssymb}

% --- Graphics ---
\usepackage{graphicx}
\usepackage{float}
\usepackage{caption}
\usepackage{subcaption}
\captionsetup{font=small, labelfont=bf}

% --- Tables ---
\usepackage{booktabs}
\usepackage{array}
\usepackage{multirow}
\usepackage{longtable}

% --- Colors ---
\usepackage{xcolor}
\definecolor{fp32color}{HTML}{0072B2}
\definecolor{fp16color}{HTML}{009E73}
\definecolor{int8ptcolor}{HTML}{E69F00}
\definecolor{int8pccolor}{HTML}{D55E00}
\definecolor{lightgray}{HTML}{F4F4F4}
\definecolor{posgreen}{HTML}{1A7A1A}
\definecolor{negred}{HTML}{C0392B}

% --- Hyperlinks ---
\usepackage[
  colorlinks=true,
  linkcolor=fp32color,
  citecolor=fp32color,
  urlcolor=fp32color
]{hyperref}

% --- Misc ---
\usepackage{microtype}
\usepackage{parskip}
\setlength{\parindent}{0pt}
\setlength{\parskip}{6pt}
\usepackage{enumitem}

% --- Header / Footer ---
\usepackage{fancyhdr}
\pagestyle{fancy}
\fancyhf{}
\lhead{\small Tudose Alexandru --- Faza 4: Scalare la Modele Mai Mari}
\rhead{\small 2026}
\cfoot{\thepage}
\renewcommand{\headrulewidth}{0.4pt}

% Helper for delta coloring
\newcommand{\dpos}[1]{\textcolor{posgreen}{\textbf{#1}}}
\newcommand{\dneg}[1]{\textcolor{negred}{#1}}
\newcommand{\dneut}[1]{#1}

% ===========================================================================
\begin{document}
% ===========================================================================

% --- Title page ---
\begin{titlepage}
  \centering
  {\large Universitatea Transilvania din Brașov\\}
  {\small Facultatea de Matematică și Informatică\\}
  {\small Specializarea: Informatică Aplicată\\[2cm]}

  \rule{\linewidth}{0.4pt}\\[0.5cm]
  {\LARGE\bfseries Faza 4: Scalare la Modele Mai Mari\\[0.3cm]}
  {\Large\bfseries ViT-Tiny / ViT-Small / ViT-Base\\[0.2cm]}
  {\Large\bfseries pe ImageNet-1k Validation\\[0.2cm]}
  {\large Evaluare pe NVIDIA A100 (Google Colab Pro+)}\\[0.3cm]
  \rule{\linewidth}{0.4pt}\\[2cm]

  \begin{minipage}{0.45\textwidth}
    \textbf{Student:}\\
    Tudose Alexandru\\
    Grupa 10LF333\\
    An III, Informatică Aplicată
  \end{minipage}
  \hfill
  \begin{minipage}{0.45\textwidth}
    \raggedleft
    \textbf{Coordonator:}\\
    Conf.\ dr.\ ing.\ Honorius Gâlmeanu\\
    Informatică Aplicată
  \end{minipage}\\[3cm]

\end{titlepage}

% --- TOC ---
\tableofcontents
\newpage

% ===========================================================================
\section{Introducere}
% ===========================================================================

Fazele~1--3 ale acestei lucrări au stabilit comportamentul cuantizării numerice
(FP16, INT8 per-tensor, INT8 per-channel) pe modelul \textbf{ViT-Tiny}
(\texttt{vit\_tiny\_patch16\_224}, 5.7M parametri) evaluat pe setul de validare
ImageNet-1k pe hardware Apple M4 (MPS). Faza~4 extinde analiza la
\textbf{modele de capacitate superioară} și la \textbf{hardware GPU dedicat},
răspunzând la întrebarea centrală:

\begin{quote}
\textit{Concluziile obținute pe ViT-Tiny/M4 se mențin la scară mai mare,
pe ViT-Small (22M) și ViT-Base (86M), rulând pe NVIDIA A100?}
\end{quote}

Toate experimentele din Faza~4 au rulat pe \textbf{Google Colab Pro+},
cu acces garantat la un GPU NVIDIA A100 (40 GB HBM2), folosind
PyTorch~2.10.0+cu128.

\subsection{Modele evaluate}

\begin{table}[H]
  \centering
  \caption{Modelele ViT evaluate în Faza~4}
  \small
  \begin{tabular}{p{5.2cm}lrrl}
    \toprule
    \textbf{Model} & \textbf{Variantă} & \textbf{Param.} &
    \textbf{Dim.} & \textbf{Sursă} \\
    \midrule
    \texttt{vit\_tiny\_patch16\_224}  & ViT-Tiny  &  5.7M & 192 &
      timm (pretrained) \\
    \texttt{vit\_small\_patch16\_224} & ViT-Small & 21.7M & 384 &
      timm (pretrained) \\
    \texttt{vit\_base\_patch16\_224}  & ViT-Base  & 86.6M & 768 &
      timm (pretrained) \\
    \bottomrule
  \end{tabular}
  \label{tab:models}
\end{table}

\subsection{Metodologie}

Protocolul experimental este identic cu cel din Fazele~1--3:
\begin{itemize}
  \item Evaluare pe toate cele 50\,000 de imagini ImageNet-1k validation;
  \item \textbf{FP32} — baseline, model \texttt{model.float()};
  \item \textbf{FP16} — \texttt{model.half()}, static;
  \item \textbf{INT8-pt} — cuantizare weight-only per-tensor, scalare liniară,
    skip LayerNorm și head de clasificare;
  \item \textbf{INT8-pc} — cuantizare weight-only per-channel (per rând).
\end{itemize}

Latențele sunt raportate ca \emph{medie per batch} (ms/batch), excluzând
primele 3 batch-uri de warmup. Dimensiunea batch-ului a fost 128 pentru
ViT-Tiny și ViT-Small și 64 pentru ViT-Base (limitare memorie VRAM pentru
FP32 la 330 MB).

% ===========================================================================
\section{Rezultate per model}
% ===========================================================================

\subsection{ViT-Tiny (5.7M parametri)}

\begin{table}[H]
  \centering
  \caption{ViT-Tiny — rezultate Faza~4 (NVIDIA A100, batch 128)}
  \begin{tabular}{lrrrr}
    \toprule
    \textbf{Format} & \textbf{Accuracy} & \textbf{$\Delta$ FP32 (pp)} &
    \textbf{Memorie (MB)} & \textbf{Latență (ms/batch)} \\
    \midrule
    \textcolor{fp32color}{\textbf{FP32}}    & 75.454\% & ---                   & 21.81 & 29.64 \\
    \textcolor{fp16color}{\textbf{FP16}}    & 75.458\% & \dpos{+0.004}         & 10.91 & \phantom{0}8.96 \\
    \textcolor{int8ptcolor}{\textbf{INT8-pt}} & 75.128\% & \dneg{$-$0.326}    & \phantom{0}6.62 & 30.69 \\
    \textcolor{int8pccolor}{\textbf{INT8-pc}} & 75.418\% & \dneg{$-$0.036}    & \phantom{0}6.70 & 30.67 \\
    \bottomrule
  \end{tabular}
  \label{tab:tiny}
\end{table}

\textbf{Observații ViT-Tiny/A100:}
FP16 obține un speedup de $\mathbf{3.31\times}$ față de FP32
(vs.\ $1.22\times$ pe M4~MPS), confirmând că A100 utilizează eficient
tensor cores. Reducerea de memorie este $2.00\times$ (exact, conform
așteptărilor teoretice). INT8-pc reduce degradarea de la $-0.326$~pp la
$-0.036$~pp, menținând pattern-ul stabilit în Faza~3.

\subsection{ViT-Small (21.7M parametri)}

\begin{table}[H]
  \centering
  \caption{ViT-Small — rezultate Faza~4 (NVIDIA A100, batch 128)}
  \begin{tabular}{lrrrr}
    \toprule
    \textbf{Format} & \textbf{Accuracy} & \textbf{$\Delta$ FP32 (pp)} &
    \textbf{Memorie (MB)} & \textbf{Latență (ms/batch)} \\
    \midrule
    \textcolor{fp32color}{\textbf{FP32}}    & 81.382\% & ---                  & 84.12  & 73.09 \\
    \textcolor{fp16color}{\textbf{FP16}}    & 81.394\% & \dpos{+0.012}        & 42.06  & 15.04 \\
    \textcolor{int8ptcolor}{\textbf{INT8-pt}} & 81.296\% & \dneg{$-$0.086}   & 23.37  & 73.77 \\
    \textcolor{int8pccolor}{\textbf{INT8-pc}} & 81.396\% & \dpos{+0.014}     & 23.53  & 73.65 \\
    \bottomrule
  \end{tabular}
  \label{tab:small}
\end{table}

\textbf{Observații ViT-Small/A100:}
FP16 atinge $\mathbf{4.86\times}$ speedup. Remarcabil: \textbf{INT8-pc
depășește FP32 cu +0.014~pp} --- primul caz din studiu în care cuantizarea
îmbunătățește marginal acuratețea față de baseline, consistent cu
observația din MLP-only INT8 pe ViT-Tiny (+0.030~pp). Degradarea INT8-pt
scade la $-0.086$~pp față de $-0.322$~pp pe ViT-Tiny.

\subsection{ViT-Base (86.6M parametri)}

\begin{table}[H]
  \centering
  \caption{ViT-Base — rezultate Faza~4 (NVIDIA A100, batch 64)}
  \begin{tabular}{lrrrr}
    \toprule
    \textbf{Format} & \textbf{Accuracy} & \textbf{$\Delta$ FP32 (pp)} &
    \textbf{Memorie (MB)} & \textbf{Latență (ms/batch)} \\
    \midrule
    \textcolor{fp32color}{\textbf{FP32}}    & 85.102\% & ---                  & 330.23 & 128.77 \\
    \textcolor{fp16color}{\textbf{FP16}}    & 85.110\% & \dpos{+0.008}        & 165.12 & \phantom{0}18.24 \\
    \textcolor{int8ptcolor}{\textbf{INT8-pt}} & 85.080\% & \dneg{$-$0.022}   & \phantom{0}87.23 & 129.78 \\
    \textcolor{int8pccolor}{\textbf{INT8-pc}} & 85.100\% & \dneg{$-$0.002}   & \phantom{0}87.55 & 129.88 \\
    \bottomrule
  \end{tabular}
  \label{tab:base}
\end{table}

\textbf{Observații ViT-Base/A100:}
FP16 atinge $\mathbf{7.06\times}$ speedup --- cel mai mare speedup din
întreg studiul. Degradarea INT8-pc coboară la o valoare aproape neglijabilă
de $-0.002$~pp, practic identică cu FP32. INT8-pt menține $-0.022$~pp,
extrem de redus față de $-0.322$~pp observat pe ViT-Tiny.

% ===========================================================================
\section{Tabel comparativ cross-model}
% ===========================================================================

Tabelul~\ref{tab:crossmodel} centralizează toate rezultatele Fazei~4
pentru o comparație directă între cele trei arhitecturi și patru formate.

\begin{table}[H]
  \centering
  \caption{Tabel comparativ cross-model --- Faza~4 (NVIDIA A100)}
  \resizebox{\linewidth}{!}{%
  \begin{tabular}{llrr>{\raggedleft\arraybackslash}p{1.8cm}>{\raggedleft\arraybackslash}p{2.2cm}}
    \toprule
    \textbf{Model} & \textbf{Format} &
    \textbf{Accuracy} & \textbf{$\Delta$ FP32 (pp)} &
    \textbf{Memorie (MB)} & \textbf{Latență (ms/batch)} \\
    \midrule
    \multirow{4}{*}{\textbf{ViT-Tiny} (5.7M)}
      & \textcolor{fp32color}{FP32}    & 75.454\% & ---           & 21.81  &  29.64 \\
      & \textcolor{fp16color}{FP16}    & 75.458\% & \dpos{+0.004} & 10.91  &   8.96 \\
      & \textcolor{int8ptcolor}{INT8-pt} & 75.128\% & \dneg{$-$0.326} & 6.62 & 30.69 \\
      & \textcolor{int8pccolor}{INT8-pc} & 75.418\% & \dneg{$-$0.036} & 6.70 & 30.67 \\
    \midrule
    \multirow{4}{*}{\textbf{ViT-Small} (22M)}
      & \textcolor{fp32color}{FP32}    & 81.382\% & ---           & 84.12  &  73.09 \\
      & \textcolor{fp16color}{FP16}    & 81.394\% & \dpos{+0.012} & 42.06  &  15.04 \\
      & \textcolor{int8ptcolor}{INT8-pt} & 81.296\% & \dneg{$-$0.086} & 23.37 & 73.77 \\
      & \textcolor{int8pccolor}{INT8-pc} & 81.396\% & \dpos{+0.014} & 23.53 & 73.65 \\
    \midrule
    \multirow{4}{*}{\textbf{ViT-Base} (86M)}
      & \textcolor{fp32color}{FP32}    & 85.102\% & ---           & 330.23 & 128.77 \\
      & \textcolor{fp16color}{FP16}    & 85.110\% & \dpos{+0.008} & 165.12 &  18.24 \\
      & \textcolor{int8ptcolor}{INT8-pt} & 85.080\% & \dneg{$-$0.022} & 87.23 & 129.78 \\
      & \textcolor{int8pccolor}{INT8-pc} & 85.100\% & \dneg{$-$0.002} & 87.55 & 129.88 \\
    \bottomrule
  \end{tabular}}%
  \label{tab:crossmodel}
\end{table}

\begin{figure}[H]
  \centering
  \includegraphics[width=\linewidth]{%
    ../results/Phase4/cross_model_comparison/00_accuracy_cross_model.png}
  \caption{Acuratețe Top-1 (\%) pentru toate cele 3 modele și 4 formate,
    evaluare pe NVIDIA~A100. Barele sunt grupate pe model; în cadrul
    fiecărui grup, ordinea este FP32, FP16, INT8-pt, INT8-pc.}
  \label{fig:accuracy-cross}
\end{figure}

\begin{figure}[H]
  \centering
  \includegraphics[width=0.85\linewidth]{%
    ../results/Phase4/cross_model_comparison/01_degradation_heatmap.png}
  \caption{Heatmap de degradare a acurateței față de FP32 (pp).
    Culorile verzi indică degradare minimă sau chiar îmbunătățire;
    culorile roșii indică degradare mai mare.
    Cel mai întunecat roșu apare la INT8-pt / ViT-Tiny ($-0.326$~pp).}
  \label{fig:heatmap}
\end{figure}

% ===========================================================================
\section{Analiză: menținerea pattern-urilor la scară mai mare}
% ===========================================================================

\subsection{Pattern 1: Per-channel vs per-tensor --- îmbunătățire consistentă}

Cuantizarea per-channel îmbunătățește semnificativ acuratețea față de
per-tensor la \emph{toate} cele trei modele:

\begin{table}[H]
  \centering
  \caption{Ameliorarea INT8-pc față de INT8-pt (pp)}
  \begin{tabular}{lrrr}
    \toprule
    \textbf{Model} & \textbf{INT8-pt $\Delta$} & \textbf{INT8-pc $\Delta$} &
    \textbf{Ameliorare} \\
    \midrule
    ViT-Tiny  (5.7M) & $-0.326$~pp & $-0.036$~pp & \dpos{$+0.290$~pp} \\
    ViT-Small (22M)  & $-0.086$~pp & $+0.014$~pp & \dpos{$+0.100$~pp} \\
    ViT-Base  (86M)  & $-0.022$~pp & $-0.002$~pp & \dpos{$+0.020$~pp} \\
    \bottomrule
  \end{tabular}
  \label{tab:ptvspc}
\end{table}

Observăm că ameliorarea absolută în pp \emph{scade} cu dimensiunea
modelului. Aceasta nu reflectă o scădere a eficacității per-channel, ci
faptul că modelele mai mari au distribuții de ponderi mai uniforme
(embed\_dim mai mare), deci și per-tensor funcționează mai bine.
Raportul dintre cele două rămâne constant: per-channel este întotdeauna
superior per-tensor.

\begin{figure}[H]
  \centering
  \includegraphics[width=0.85\linewidth]{%
    ../results/Phase4/cross_model_comparison/06_int8_degradation_scaling.png}
  \caption{Evoluția degradării INT8-pt și INT8-pc pe măsura creșterii
    dimensiunii modelului. Ambele curbe converg spre zero, dar per-channel
    (portocaliu) rămâne constant aproape de linia FP32.}
  \label{fig:int8-scaling}
\end{figure}

\subsection{Pattern 2: FP16 speedup crește odată cu dimensiunea modelului}

Acesta este cel mai semnificativ fenomen observat în Faza~4, specific
hardware-ului GPU:

\begin{table}[H]
  \centering
  \caption{FP16 speedup față de FP32 (NVIDIA A100)}
  \begin{tabular}{lrrrr}
    \toprule
    \textbf{Model} & \textbf{FP32 lat.} & \textbf{FP16 lat.} &
    \textbf{Speedup A100} & \textbf{Speedup M4 MPS} \\
    \midrule
    ViT-Tiny  (5.7M) &  29.64 ms &  8.96 ms & $\mathbf{3.31\times}$ & $1.22\times$ \\
    ViT-Small (22M)  &  73.09 ms & 15.04 ms & $\mathbf{4.86\times}$ & --- \\
    ViT-Base  (86M)  & 128.77 ms & 18.24 ms & $\mathbf{7.06\times}$ & --- \\
    \bottomrule
  \end{tabular}
  \label{tab:fp16speedup}
\end{table}

Pe M4~MPS (Faza~1), FP16 oferea un modest $1.22\times$ speedup.
Pe A100, același ViT-Tiny obține $3.31\times$, iar la ViT-Base ajunge
la $7.06\times$. Explicația constă în natura calculului:
\begin{itemize}
  \item \textbf{Modele mici, batch mic} (ViT-Tiny, bs=128): calculul este
    parțial \emph{memory-bound}; transferul de date limitează speedup-ul.
  \item \textbf{Modele mari} (ViT-Base, bs=64): calculul devine
    \emph{compute-bound}; tensor cores A100 (BF16/FP16) lucrează la
    aproape $2\times$ flops față de FP32, cu overhead de transfer
    proporțional mai mic.
\end{itemize}

\begin{figure}[H]
  \centering
  \includegraphics[width=0.78\linewidth]{%
    ../results/Phase4/cross_model_comparison/05_fp16_speedup_scaling.png}
  \caption{Speedup FP16 față de FP32 pe NVIDIA A100 în funcție de
    dimensiunea modelului. Creșterea este sub-liniară dar consistentă,
    de la $3.31\times$ (Tiny) la $7.06\times$ (Base).}
  \label{fig:fp16-speedup}
\end{figure}

\subsection{Pattern 3: Reducerea de memorie --- exact 2× pentru FP16, $\approx$3.3--3.8× pentru INT8}

\begin{table}[H]
  \centering
  \caption{Raport de reducere memorie față de FP32}
  \begin{tabular}{lrrrr}
    \toprule
    \textbf{Model} & \textbf{FP32 (MB)} & \textbf{FP16} &
    \textbf{INT8-pt} & \textbf{INT8-pc} \\
    \midrule
    ViT-Tiny  (5.7M)  &   21.81 & $2.00\times$ & $3.29\times$ & $3.25\times$ \\
    ViT-Small (22M)   &   84.12 & $2.00\times$ & $3.60\times$ & $3.57\times$ \\
    ViT-Base  (86M)   &  330.23 & $2.00\times$ & $3.79\times$ & $3.77\times$ \\
    \bottomrule
  \end{tabular}
  \label{tab:memory}
\end{table}

Reducerea FP16 este \emph{exact} $2.00\times$ indiferent de model ---
demonstrație directă a njumătățirii lățimii de biți. Reducerea INT8
\emph{crește} ușor cu dimensiunea modelului (de la $3.29\times$ la
$3.79\times$), deoarece modelele mai mari au o proporție mai mică din
parametri stocați ca bias-uri și scale-uri LayerNorm (care rămân în FP32).

\begin{figure}[H]
  \centering
  \includegraphics[width=0.85\linewidth]{%
    ../results/Phase4/cross_model_comparison/02_memory_scaling.png}
  \caption{Evoluția memoriei (MB) pe măsura creșterii modelului.
    Rețineți că axele sunt diferite între modele --- graficul evidențiază
    divergența între formatele FP32/FP16 și INT8, care se accentuează
    la ViT-Base (330 MB FP32 vs 87 MB INT8).}
  \label{fig:memory-scaling}
\end{figure}

\subsection{Pattern 4: INT8 fără speedup pe GPU --- confirmat la toate modelele}

La toate cele trei modele, INT8 nu produce speedup față de FP32 pe A100:
latența este practic identică ($0.97$--$0.99\times$). Implementarea noastră
\texttt{QuantizedLinear} stochează ponderile în INT8 dar \emph{dequantizează
la FP32 înainte de \texttt{torch.nn.functional.linear}}. Prin urmare,
avantajul de memorie se menține, dar nu există kernel INT8 nativ care să
reducă latența. Acest comportament este \emph{independent de hardware} ---
aceeași limita a fost observată și pe M4~MPS.

% ===========================================================================
\section{Comparație hardware: Apple M4 MPS vs NVIDIA A100}
% ===========================================================================

Pentru ViT-Tiny, dispunem de date pe ambele platforme, permițând o comparație
directă. Notă: dimensiunile de batch diferă (M4: 64, A100: 128), deci
valorile sunt normalizate la \emph{ms per imagine}.

\begin{table}[H]
  \centering
  \caption{Comparație hardware ViT-Tiny --- latență per imagine (ms/imagine)}
  \begin{tabular}{lrrrr}
    \toprule
    \textbf{Format} & \textbf{M4 MPS} & \textbf{A100 GPU} &
    \textbf{Speedup A100/M4} \\
    \midrule
    FP32    & 2.578 ms & 0.232 ms & $\mathbf{11.1\times}$ \\
    FP16    & 2.111 ms & 0.070 ms & $\mathbf{30.2\times}$ \\
    INT8-pt & 4.299 ms & 0.240 ms & $\mathbf{17.9\times}$ \\
    INT8-pc & 4.545 ms & 0.240 ms & $\mathbf{18.9\times}$ \\
    \bottomrule
  \end{tabular}
  \label{tab:hardware}
\end{table}

\textbf{Observații cheie:}
\begin{enumerate}
  \item \textbf{FP32}: A100 este de $11.1\times$ mai rapid --- beneficiaza
    de paralelism masiv și bandwidth HBM2 ($\approx$2~TB/s vs $\approx$100~GB/s MPS).
  \item \textbf{FP16}: speedup A100/M4 sare la $30.2\times$, deoarece
    A100 utilizează tensor cores pentru FP16 (eficiență $\approx 2\times$
    față de FP32), în timp ce M4~MPS are un speedup modest de $1.22\times$.
  \item \textbf{INT8}: pe M4, INT8 este \emph{mai lent} decât FP32 (overhead
    de dequantizare software); pe A100 latența este similară cu FP32.
    Speedup-ul A100/M4 ajunge la $17.9$--$18.9\times$ deoarece startul
    pe M4 este penalizat și mai mult.
\end{enumerate}

\begin{figure}[H]
  \centering
  \includegraphics[width=\linewidth]{%
    ../results/Phase4/cross_model_comparison/04_hardware_comparison.png}
  \caption{Latență per imagine (ms/imagine) pentru ViT-Tiny pe Apple~M4~MPS
    (Fazele~1--3) și NVIDIA~A100 (Faza~4). A100 este consistent mai rapid,
    cu cel mai mare avantaj la FP16 ($30.2\times$).}
  \label{fig:hardware}
\end{figure}

\textbf{Concluzie hardware:} FP16 este formatul care beneficiază cel mai mult
de pe urma hardware-ului GPU dedicat. Pe M4~MPS, FP16 oferă doar $1.22\times$
speedup din cauza suportului limitat pentru tensor operations în half-precision;
pe A100, același format devine de departe cel mai rapid, cu speedup de
$3.31\times$ (Tiny) până la $7.06\times$ (Base).

% ===========================================================================
\section{Acuratețe vs. dimensiunea modelului}
% ===========================================================================

Un rezultat important al Fazei~4 este că acuratețea \emph{crește}
semnificativ cu dimensiunea modelului, iar impactul cuantizării
\emph{scade} în valoare absolută:

\begin{figure}[H]
  \centering
  \includegraphics[width=\linewidth]{%
    ../results/Phase4/cross_model_comparison/03_latency_scaling.png}
  \caption{Latență per batch (ms/batch) în funcție de dimensiunea modelului
    și formatul de cuantizare. FP16 scalează aproape liniar, în timp ce
    FP32 și INT8 au o creștere mai accentuată. La ViT-Base, FP16 rămâne
    sub 20~ms/batch față de 128~ms FP32.}
  \label{fig:latency-scaling}
\end{figure}

Acuratețele baseline FP32 sunt:
ViT-Tiny $75.45\%$ $\to$ ViT-Small $81.38\%$ $\to$ ViT-Base $85.10\%$.
Aceasta reflectă capacitatea crescută a modelelor mai mari de a reprezenta
caracteristicile vizuale din ImageNet-1k.

Degradarea maximă din cuantizare (INT8-pt) scade dramatic cu dimensiunea:
$-0.326$~pp (Tiny) $\to$ $-0.086$~pp (Small) $\to$ $-0.022$~pp (Base).
Modelele mai mari sunt \textbf{mai robuste la cuantizare}, deoarece au
distribuții de ponderi mai netede (spread mai mic pe dimensiunea embed\_dim
mai mare), permițând un scalar de cuantizare mai precis.

% ===========================================================================
\section{Concluzii și ghid practic de cuantizare}
% ===========================================================================

\subsection{Concluzii sintetice Faza~4}

\begin{enumerate}
  \item \textbf{Pattern-urile se mențin la scară mai mare.}
    Superioritatea INT8-pc față de INT8-pt, comportamentul FP16
    și reducerile de memorie sunt consistente de la ViT-Tiny la ViT-Base.

  \item \textbf{Modele mai mari = cuantizare mai sigură.}
    La ViT-Base, INT8-pc pierde doar $0.002$~pp față de FP32 ---
    practic neglijabil. La ViT-Tiny, aceeași metodă pierdea $0.036$~pp.

  \item \textbf{FP16 devine din ce în ce mai atractiv pe GPU.}
    Speedup-ul de $7\times$ la ViT-Base, cu memorie exact $2\times$
    mai mică și fără degradare de acuratețe, face din FP16 alegerea
    evidentă pentru inference pe GPU dedicat.

  \item \textbf{INT8 fără kernel nativ nu aduce speedup pe GPU.}
    Implementarea weight-only cu dequantizare la FP32 este utilă
    exclusiv pentru reducerea memoriei, nu pentru latență.
\end{enumerate}

\subsection{Ghid practic de cuantizare --- sinteza tuturor fazelor}

Pe baza rezultatelor din Fazele~1--4 și a experimentelor adiționale
(MLP-only INT8, FP8 E4M3FN), propunem următorul ghid de decizie:

\begin{table}[H]
  \centering
  \caption{Ghid practic --- alegerea formatului de cuantizare}
  \small
  \begin{tabular}{p{3.5cm}p{3.5cm}p{4cm}l}
    \toprule
    \textbf{Constrângere} & \textbf{Hardware} &
    \textbf{Recomandare} & \textbf{Pierdere acc.} \\
    \midrule
    Inference rapidă, & GPU (A100/V100/ & \textbf{FP16 static} &
      $\approx 0$ \\
    memorie $2\times$ & RTX 3090+) & (\texttt{model.half()}) & \\
    \addlinespace
    Memorie $3\times$, & Orice & \textbf{INT8 per-channel} &
      $< 0.04$~pp \\
    acuratețe înaltă & (GPU/CPU/MPS) & (skip LayerNorm, head) & \\
    \addlinespace
    Memorie $3\times$, & Orice & \textbf{INT8 per-tensor} &
      $\approx 0.1$--$0.3$~pp \\
    simplitate & & & \\
    \addlinespace
    Memorie redusă, & GPU cu suport & \textbf{FP8 E4M3FN} &
      $0.1$--$0.2$~pp \\
    viteză FP16 & FP8 nativ & (per-tensor > INT8-pt) & \\
    \addlinespace
    Acuratețe $\geq$ FP32, & Orice & \textbf{INT8-pc pe MLP} &
      $+0.03$~pp \\
    memorie moderată & & (skip atenție) & \\
    \addlinespace
    Deployment fără & Orice & \textbf{FP32} &
      --- \\
    compromis & & & \\
    \bottomrule
  \end{tabular}
  \label{tab:guide}
\end{table}

\subsection{Ierarhia formatelor --- concluzie generală}

Din perspectiva acurateței (cel mai important criteriu în producție):
{\small\[
  \underbrace{\text{INT8-mlp-pc}}_{\text{+0.030 pp}} \geq
  \underbrace{\text{FP16}}_{\approx 0} \approx
  \underbrace{\text{INT8-pc}}_{-0.002\text{--}-0.036} >
  \underbrace{\text{FP8-pc}}_{-0.100} >
  \underbrace{\text{FP8-pt}}_{-0.198} >
  \underbrace{\text{INT8-pt}}_{-0.022\text{--}-0.326}
\]}

Din perspectiva memoriei:
\[
  \text{FP32} > \text{FP16}\ (2\times) > \text{INT8/FP8}\ (3.3\text{--}3.8\times)
\]

Din perspectiva latenței pe GPU (A100):
\[
  \text{FP16}\ (7\times \text{ speedup}) \gg
  \text{FP32} \approx \text{INT8}\ (\text{fără speedup, impl. noastră})
\]

\textbf{Concluzie finală:} Cuantizarea numerică a Vision Transformers este
matură și robustă la toate scările analizate. FP16 este alegerea optimă
pentru GPU modere, oferind viteza maximă fără pierdere de acuratețe.
INT8 per-channel reprezintă compromisul ideal pentru deployment cu
constrângeri de memorie. Modelele mai mari (ViT-Base, 86M) sunt cel mai
puțin afectate de cuantizare, cu degradări sub $0.02$~pp.

% ===========================================================================
\section{Detalii tehnice și reproducibilitate}
% ===========================================================================

\begin{table}[H]
  \centering
  \caption{Configurație hardware și software --- Faza~4}
  \small
  \begin{tabular}{lp{9.5cm}}
    \toprule
    \textbf{Parametru} & \textbf{Valoare} \\
    \midrule
    Hardware & NVIDIA A100 (40 GB HBM2), Google Colab Pro+ \\
    PyTorch & 2.10.0+cu128 \\
    CUDA & 12.8 \\
    timm & pretrained checkpoints ImageNet-1k \\
    Dataset & ImageNet-1k validation, 50\,000 imagini \\
    Batch size & 128 (Tiny/Small), 64 (Base) \\
    Warmup & 3 batch-uri excluse din timing \\
    Latență & medie pe toate batch-urile post-warmup \\
    \midrule
    Script principal & \texttt{scripts/phase4/evaluate\_model.py} \\
    Script comparație & \texttt{scripts/phase4/compare\_cross\_model.py} \\
    Rezultate & \texttt{results/Phase4/\allowbreak<model>/\allowbreak metrics/results.json} \\
    \bottomrule
  \end{tabular}
  \label{tab:config}
\end{table}

Toate rezultatele sunt stocate în fișiere JSON și pot fi reproduse rulând
scriptul de evaluare cu argumentele \texttt{-{}-model} și
\texttt{-{}-batch-size} pe orice hardware cu CUDA disponibil.

% ===========================================================================
\end{document}
